\paragraph{The compact real factor and the definite lattice mass.}
Put \(n=[E:\mathbb Q]=2f\), and write \(O=O_E\). The two numerical
factors in the preceding comparison have the following bounds.
\begin{proposition}\label{bergvol:compactmass}
For the real density specified below and the genus of the standard
positive definite Hermitian \(O\)-lattice of rank \(r\ge1\),
\[
 c_r=\operatorname{vol}(U(r))
     =\prod_{j=1}^r\frac{2\pi^j}{(j-1)!}\le13^{r^2},
 \qquad
 \operatorname{mass}_I
     \le [3n^2r^2R_E^2]^{2nr^2}
     \le[3n^2r^4R_E^2]^{3nr^2}.
\]
\end{proposition}
\begin{proof}
In the skew-Hermitian Lie algebra use the density \(dx\) given by
the imaginary diagonal entries and the real and imaginary parts of
the entries strictly below the diagonal. To compare it with the real
measure of the tube calculation, let \(\lambda\) be the additive
measure for which \((X,Y)\mapsto X+iY\) carries
\(\lambda\otimes\lambda\) to entrywise Lebesgue measure on
\(M_r(\mathbb C)\). Each diagonal block has absolute real determinant
one. Each off-diagonal pair gives the two complex entries
\(z+iw,-\bar z-i\bar w\), a block of absolute real determinant four.
Hence \(\lambda=2^{\binom r2}dx\). The specified real Haar density
at the identity, in exponential coordinates, is
\(2^{-\binom r2}\lambda=dx\).
\src{src/Formalism/LocalRingCardinality/DataCD.lean}{42}{Complexification measure on the real Lie algebra};
\src{src/Formalism/LocalToGlobal/BesselHermitianTheta/MinfDensity.lean}{13}{Real density relative to complexification}.

The map \(U(j)\to S^{2j-1}\), \(g\mapsto ge_1\), has fibre
\(U(j-1)\). At the identity its differential selects the first
column, whose independent coordinates give Euclidean tangent
measure on the sphere at \(e_1\). The remaining coordinates give
exactly the chosen density on \(U(j-1)\). By invariance the same
quotient measure is Euclidean surface measure on the whole sphere.
Starting with \(c_0=1\) gives the displayed product. For \(j\le r\)
each factor is at most \(2\cdot4^r\le8^r\).
In particular compact probability Haar has density
\(2^{-\binom r2}/c_r\) relative to \(\lambda\). For a real form
\(U(p,q)\), \(p+q=r\), the compact-dual convention prescribes this
same scalar relative to its own complexification measure. Thus
multiplication by \(c_r\) converts that convention to the specified
real density at each place.
\src{src/Formalism/LocalRingCardinality/DataCD.lean}{57}{Matching of the real density across signatures};
\src{src/ReciprocityTheory/ReductiveGroups/DiscreteSeries/FormalDegree/VolUFactProdMulLe.lean}{14}{Compact real volume factor}.

For the second assertion, every lattice \(L\) in the standard genus
is Hermitian self-dual. Indeed this holds locally for \(O^r\), and
localization commutes with the dual of a projective lattice. One may
use the conjugation-stable semilocal rings at pairs
\(\{\mathfrak p,\bar{\mathfrak p}\}\) here, including at split primes.
We bound the number of isometry classes of all totally positive
definite self-dual Hermitian \(O\)-lattices of rank \(r\).
Give the underlying real space the scalar product
\(\tfrac12\operatorname{Tr}_{E/\mathbb Q}h(x,y)\). Its squared norm
is \(\sum_\tau h_\tau(x_\tau,x_\tau)\), with one summand per conjugate
pair. For the absolute different \(\mathfrak D\) the Euclidean
dual lattice is \(L^*=2\mathfrak D^{-1}L\). Indeed testing the
pairing on all \(O\)-multiples of each lattice vector gives precisely
the codifferent condition, and \(L\) is Hermitian self-dual.
For \(0\ne y\in\mathfrak D^{-1}L\) we have
\(h(y,y)\in\mathfrak D^{-1}\bar{\mathfrak D}^{-1}\), whence
\(\lvert N_{E/\mathbb Q}h(y,y)\rvert\ge\Delta^{-2}\).
The arithmetic-geometric mean inequality gives
\(\|y\|^2\ge f\Delta^{-2/n}=fR_E^{-2}\).
Thus \(\lambda_1(L^*)\ge 2\sqrt f/R_E\).
By Banaszczyk's Euclidean transference theorem in dimension \(nr\),
\(\lambda_{nr}(L)\lambda_1(L^*)\le nr\)
\cite[Theorem~2.1]{banaszczyk}. Among real-linearly independent
vectors realizing this last minimum we may choose an \(E\)-basis
\(x_1,\ldots,x_r\) with \(\|x_i\|\le T=nrR_E\).

Write \(G=(h(x_i,x_j))\) and \(L_0=\sum O x_i\). The entries of
\(G\) are integral and have modulus at most \(T^2\) at every
embedding. There are at most \((1+2T^2)^n\) choices per entry.
Indeed the points of \(O\) in its Minkowski space have mutual sup distance
at least one by the norm inequality, so open polydiscs of radius
\(1/2\) around these points are disjoint. For a fixed \(G\) the marked
rational Hermitian space and \(L_0\) are fixed up to isometry.
The inclusions \(L_0\subseteq L\subseteq L_0^\#\), where \(\#\)
denotes the Hermitian dual, locate \(L/L_0\) in a finite module
\(A=L_0^\#/L_0\) of order \(\lvert N_{E/\mathbb Q}\det G\rvert
\le T^{2nr}\). The last inequality follows by Gram's determinant
inequality at every embedding.

Each \(O\)-submodule of \(A\) has an ordered generating tuple of
length \(r\). Locally at each prime in the support its inverse image
in \((L_0^\#)_{\mathfrak p}\) is free of rank \(r\) over a discrete
valuation ring, so its image has \(r\) generators. Combine these
tuples coordinate by coordinate under the primary decomposition of
the finite module. Consequently there are at most \(|A|^r\)
intermediate lattices for the fixed Gram matrix. All isometry
classes together number at most
\[
 (1+2T^2)^{nr^2}T^{2nr^2}\le (3T^2)^{2nr^2}.
\]
The definite standard mass is the sum of
\(1/|\operatorname{Aut}_O(L,h)|\) over the classes in its genus.
Each group is finite, since it is a discrete lattice stabilizer
inside a compact real group. The class count therefore proves
the required bound, including for nonfree \(L\).
\src{src/Formalism/SharpRestriction/GlobalSchurBounds/GeneralLinearActions/GapAtLeast.lean}{20}{Definition of the definite standard mass};
\src{src/Formalism/HermitianLattices/StatementMass.lean}{11}{Reciprocal stabilizer weights}.
\end{proof}

These estimates give \(V\le V_*\) once the displayed adelic tube
inequality holds. Applying that inequality to a different signature
still requires the comparison of the adelic covolumes with their
finite lattice stabilizers and the stated real densities.
